\documentclass[11pt]{article}
\usepackage[a4paper,margin=1in]{geometry}
\usepackage{amsmath,amssymb,amsthm,booktabs,hyperref}
\newtheorem{theorem}{Theorem}
\newtheorem{lemma}{Lemma}
\title{Symmetry of four-digit continued-fraction permutiples}
\author{Principia Math In Progress}
\date{28 September 2026}
\begin{document}
\maketitle
\noindent MathDB 335304; Holt, arXiv:1702.04760v1, Conjecture 17.

\begin{theorem}[Task 1]
Let $a,b,c,d$ be positive integers, and let $(e,f,g,h)$ be a
permutation of $(a,b,c,d)$. Suppose that $d,h\ge2$ and
\[
 [a;b,c,d]=k[e;f,g,h],\qquad k\in\mathbb Z,\quad k\ge2.
\]
Then $ad=eh$ and $bc=fg$. Thus every canonical four-digit
continued-fraction permutiple is symmetric in Holt's sense.
\end{theorem}

\section*{Continuants and two reductions}
Put
\[
 P=abcd+ab+ad+cd+1,\qquad Q=bcd+b+d,
\]
and define $P',Q'$ using $(e,f,g,h)$ in the same order.
The continued fractions are $P/Q$ and $P'/Q'$, in lowest terms.
Taking integer parts gives
\begin{equation}\label{floor}
 a=ke+j,\qquad 0\le j\le k-1.
\end{equation}
In particular $a>e$ and $k\le a/e\le a$.

\begin{lemma}\label{ratio}
For any two arrangements of four positive integers,
their four-term continuants have ratio strictly less than two.
\end{lemma}
\begin{proof}
Write the sorted entries as $1\le x\le y\le z\le w$ and set
$V=xyzw$. The continuant of an arrangement is $V+S+1$, where
$S$ is the sum of the three edge products along a path through
the four entries: for $(a,b,c,d)$ the path is $b,a,d,c$.
Root this path at $x$. Each other vertex $v$ contributes an edge
of weight at least $xv$, so $S\ge x(y+z+w)$.
Also $S\le w(y+z)+yz$. Indeed, if $w$ has degree two its incident
edges sum to at most $w(y+z)$ and the other edge is at most $yz$.
If it has degree one its incident edge is at most $wz$ and the
other two sum to at most $yz+xz\le yz+wy$.
With $L_0=V+x(y+z+w)+1$ and $U_0=V+w(y+z)+yz+1$, direct expansion gives
\begin{align*}
2L_0-U_0={}&(x-1)\{yzw+2(y+z+w)\}\\
 &+(y-1)(z-1)(w-1)+y+z+w+2>0.
\end{align*}
Every such continuant therefore lies between $L_0$ and $U_0<2L_0$.
\end{proof}

Since $PQ'=kP'Q$ and $\gcd(P',Q')=1$, we have $P'\mid P$.
Lemma~\ref{ratio} now gives
\begin{equation}\label{equal}
 P=P',\qquad Q'=kQ.
\end{equation}
For completeness, set $p_2=abc+a+c$, $q_2=bc+1$ and define
$p'_2,q'_2$ analogously. The determinant identities are
\[
 Pq_2-Qp_2=1,\qquad Pq'_2-Q'p'_2=1.
\]
Subtracting them and using \eqref{equal} yields
\begin{equation}\label{divis}
 Q\mid(bc-fg),\qquad P\mid(p_2-kp'_2).
\end{equation}

\begin{lemma}\label{last}
If $h\ge k$, then $ad=eh$ and $bc=fg$.
\end{lemma}
\begin{proof}
We have $0<p_2<P$ and $P=hp'_2+ef+1>hp'_2$, so
$0<kp'_2<P$ when $h\ge k$. By \eqref{divis}, their difference,
which lies strictly between $-P$ and $P$, is zero.
The subtracted determinant identities then imply $q_2=q'_2$,
or $bc=fg$. The equality $abcd=efgh$ gives $ad=eh$.
\end{proof}

\section*{Exhausting the permutations}
The six positional permutations with $e=a$ are excluded by $a>e$.
The six with $h=a$ satisfy Lemma~\ref{last}, since $a\ge k$.
The remaining twelve are in the following table. Set
$L=b-c+d$ and $N=bc+bd-cd$. Every expression in the middle column
must vanish by \eqref{equal}.
\begin{center}
\begin{tabular}{lll}
\toprule
$(e,f,g,h)$ & $P-P'$ & Resolution\\
\midrule
$(b,a,c,d)$ & $d(a-b)$ & $a=e$, impossible\\
$(b,a,d,c)$ & $ad-bc$ & symmetry directly\\
$(b,c,a,d)$ & $b(a-c)-d(b-c)$ & Case A\\
$(b,d,a,c)$ & $aL-N$ & Case B\\
$(c,a,b,d)$ & $aL-bd$ & Case C\\
$(c,a,d,b)$ & $aL-N$ & Case D\\
$(c,b,a,d)$ & $b(a-c)$ & $a=e$, impossible\\
$(c,d,a,b)$ & $ad-bc$ & symmetry directly\\
$(d,a,b,c)$ & $b(a-c)$ & $h=a$, Lemma~\ref{last}\\
$(d,a,c,b)$ & $b(a-c)-d(b-c)$ & Case E\\
$(d,b,a,c)$ & $aL-bd$ & Case C\\
$(d,c,a,b)$ & $d(a-b)$ & $h=a$, Lemma~\ref{last}\\
\bottomrule
\end{tabular}
\end{center}
This is a list of positional permutations, so repeated digit values
cause no exception to its exhaustiveness.

\paragraph{Case C.}
Here $aL=bd>0$, so the integer $L$ is at least one. Since
$c=b+d-L$,
\[
 cL-b=(c-1)(L-1)+d-1\ge0.
\]
Thus $a=bd/L\le cd=eh$. As $k\le a/e$, it follows that
$h\ge k$ and Lemma~\ref{last} applies.

\paragraph{A common bound for Cases B and D.}
If $aL=N$ and $L\ge1$, then $a\le bc$.
For $b\le c$, this follows from $N=bc-d(c-b)\le bc$.
For $b>c$, use $L\ge d+1$ and $b-c<bc$ to obtain
$N=bc+d(b-c)<(d+1)bc\le Lbc$.
In both cases $eh=bc$, so $h\ge k$ and Lemma~\ref{last} applies.

\paragraph{Case D.}
Write $a=kc+j$ with $j\ge0$. The equality $Q'=kQ$ is
\[
 kL=j(bd+1)+b>0.
\]
Consequently $L\ge1$ and the common bound proves symmetry.

\paragraph{Case B.}
The common bound handles $L\ge1$. If $L=0$, then $c=b+d$
and $N=b^2+bd-d^2=0$, which would give
$(2d-b)^2=5b^2$, impossible for positive integers.
Suppose $L<0$ and put $u=-L\ge1$, so $c>b+d>d$.
Write $a=kb+j$ with $j\ge0$. Using $Q'=acd+c+d=kQ$
and $aL=N$ gives
\begin{equation}\label{hard}
 jQ=c(a-d),\qquad ad-bc=(a-d)(c-b),\qquad
 a-d=\frac{d^2-bc}{u}.
\end{equation}
The first identity implies $a\ge d$. If $a=d$, the second proves
$ad=bc$ and hence symmetry. Otherwise $0<a-d<d^2$.
By \eqref{divis}, $Q\mid(ad-bc)$. The first two identities in
\eqref{hard} then imply $Q\mid b(a-d)$ as well as $Q\mid c(a-d)$.
Writing $s=\gcd(b,c)$, B\'ezout's identity gives $Q\mid s(a-d)$.
But
\[
 0<s(a-d)<sd^2\le bd^2<bcd<Q,
\]
a contradiction. This completes Case B.

\paragraph{Case A.}
The numerator identity is $b(a-c)=d(b-c)$, so
\[
 a=c+d-\frac{cd}{b},\qquad b(a-b)=(b-c)(d-b).
\]
Since $a>b$, the possibility $d=b$ is excluded.
If $d>b$, then $a<d=h$, and Lemma~\ref{last} applies.
If $d<b$, the second identity forces $c>b$, and then $a<c$.
Write $a=kb+j$ with $j\ge0$. The denominator identity becomes
\[
 c(jd+1)=kb+(k-1)d.
\]
If $j\ge1$, then, because $c>a\ge kb$,
\[
 c(jd+1)\ge c(d+1)>kb(d+1)\ge k(b+d)>kb+(k-1)d,
\]
which is impossible. If $j=0$, then $a=kb$ and
$c=kb+(k-1)d$. Substituting into $b(a-c)-d(b-c)=0$
gives $(k-1)d^2=0$, again impossible.

\paragraph{Case E.}
Again $a=c+d-cd/b<c+d$. Write $a=kd+j$, $j\ge0$.
The denominator identity is
\[
 j(bc+1)=(k-1)b.
\]
It forces $j\ge1$ and hence $c<k-1$. But then
\[
 a<c+d<k+d-1\le kd\le a,
\]
a contradiction. All permutations have been resolved, proving the theorem.

\section*{Scope and reproducibility}
The theorem answers Conjecture 17 for all positive digits and every
integer multiplier $k>1$, using the source's canonical conventions.
No distinctness or bounded-digit assumption is made. It does not assert
the longer-length conjectures. The proof uses only elementary
continuant identities; no other structural theorem in the source is needed.
The accompanying verifier checks the polynomial identities, the full
permutation table, and bounded exact instances. The latter are checks,
not a substitute for the proof. This is author-agent self-review,
not independent peer review.

\section*{Proposed follow-up tasks}
These are research directions, with no solutions claimed here.
\begin{enumerate}
\setcounter{enumi}{1}
\item Determine the sharp supremum of the ratio of two four-term
continuants obtained by permuting one positive-integer multiset.
\item For each fixed integer $k>1$, parametrize all canonical
four-digit continued-fraction permutiples without repetitions.
\item Decide whether every canonical five-digit continued-fraction
permutiple preserves its numerator continuant.
\item Classify the canonical five-digit permutiples whose two
outer-pair products and middle digit are preserved by the permutation.
\end{enumerate}
\begin{thebibliography}{9}
\raggedright
\bibitem{Holt} B. V. Holt, \emph{Digit Preserving Multiplication in
Continued Fraction Representations}, arXiv:1702.04760v1 (2017),
Definitions 1 and 13, Conjecture 17.
\url{https://arxiv.org/abs/1702.04760}.
\end{thebibliography}
\end{document}
